% Bounded original-study contribution to planned P19 227.7--227.8.
\chapter{Aging, goals and supported reasoning}
\label{sec:p19-aging-goals-and-reasoning}

An age difference in the speed of processing information is not the same
result as an age difference in the quality of a choice. Neither determines
how much information a person chooses to examine, how well specialized
knowledge supports a familiar task, or whether a practiced skill can improve.
The four original records below separate these questions: a relevance and
accountability experiment, a decision-structure experiment, a longitudinal
study of skilled simulator performance, and a randomized reasoning-training
trial \cite{HessQueenEnnis2013Search,HessQueenPatterson2012Deliberation,
TaylorKennedyNodaYesavage2007Pilot,RebokEtAl2014ACTIVE}.
They do not measure a single lifespan rationality score.

The distinction between producing an inference and explaining its adequacy
in \cref{sec:p19-rationality-15-1} remains important in later adulthood.
Here the criteria additionally include information-search behavior,
experimenter-defined choice quality, domain-specific performance and scores
on serial-pattern tests. The designs and uncertainty distinctions of
\cref{sec:p19-evidence-methods} apply separately to each outcome.

\begin{hbdefinition}{def:p19-aging-comparisons}{Age differences, change, training and compensation}
A \emph{cross-sectional age difference} compares different people at
different ages. It may include differences in education, cohort, health,
recruitment and previous experience. \emph{Longitudinal change} compares
repeated observations of the same people, but can also include practice,
selective retention and changes in their environment. A randomized
\emph{training effect} compares assigned interventions under an identified
outcome and follow-up; it is not established by an association between
experience and performance.

\emph{Cognitive compensation} is the proposed maintenance of performance
through a different resource, strategy or support despite a limitation in
another component. Demonstrating good performance alone does not identify
that mechanism. In choice research, a \emph{compensatory decision strategy}
instead permits an advantage on one attribute to offset a disadvantage on
another. These two uses of compensation must not be conflated. A
noncompensatory choice procedure can still involve extensive, deliberate
search. In the study records, \(N\) denotes a participant or specified
outcome count, and \(p\) a reported statistical-test probability, not a
probability that an aging theory is true.
\end{hbdefinition}

\section{Personal relevance changes engagement, not a fixed capacity}
\label{sec:p19-aging-goals-search}

Hess, Queen and Ennis recruited 79 community-dwelling younger adults,
ages 21--41, and 81 older adults, ages 64--90, from a participant database.
Mean ages were 32.3 and 73.7. Mean education was 15.7 and 16.6 years.
Older participants scored lower on the digit-symbol and letter-number
tests, but not on vocabulary. These are relatively educated volunteer
groups, not random samples of all younger and older adults
\cite[pp.~705--706, table~1]{HessQueenEnnis2013Search}.

Each adult searched two arrays with eight alternatives and six attributes:
one for a wireless phone plan and one for a prescription drug plan. The
plans are historical experimental materials, not current purchasing or
health advice. A mouse click exposed one cell, and opening another hid the
previous cell. This makes the amount and sequence of search observable,
while imposing a memory demand absent from a display that keeps every
value visible. A handout explained the attributes and a practice array
familiarized participants with the interface. Task order was counterbalanced.
Within each age group, participants were randomly assigned either to
expect an oral justification of their decision or to receive no such
accountability instruction. Both tasks were completed by the same person.

The selective-engagement hypothesis predicts that people allocate effort
in relation to a task's self-relevance, with older adults particularly
sensitive to that relation. This is not the hypothesis that older people
cannot perform the search. The relevance manipulation was also not a
random assignment of a lifetime goal: the two plan domains differed in
content, familiarity and perceived relevance. Older adults rated the drug
plan as more relevant than younger adults did, and younger adults rated
the phone plan as more relevant. Nonetheless, most people in both age
groups had phone-plan experience; drug-plan experience was reported by
43 percent of older and 11 percent of younger participants. Familiarity,
knowledge and relevance therefore cannot be treated as identical variables
\cite[pp.~705--706]{HessQueenEnnis2013Search}.

Older adults spent longer searching overall, 239 versus 166 seconds after
one outlier in each age group was excluded. Accountability increased search
time overall. For the phone-plan task, older adults' mean search times
were 235.4 seconds under high accountability and 159.1 under low
accountability; younger adults' corresponding means were 138.4 and 132.4.
The reported age-by-accountability interaction for that task had
\(p=.02\). This is a difference in \emph{engagement time}, not proof of
better choices. For the proportion of information sampled, the overall
age-by-accountability interaction only approached the paper's criterion,
\(p=.10\), although its older-group follow-up was significant. A
significant contrast in one age group and a nonsignificant contrast in
another do not by themselves establish an interaction
\cite[pp.~706--707]{HessQueenEnnis2013Search}.

The initial process-tracing algorithm classified most searches as
satisficing. The authors then split these searches at the within-task
median amount of information sampled: extensive searches became the
noncompensatory group, while limited searches retained the satisficing
label. The thresholds were 90 percent for the drug-plan task and
77 percent for the phone-plan task. Consequently, the categories are
partly constructed from the same search quantities later used to
characterize them; those characterizations are not an independent
validation of a uniquely identified mental algorithm.

For older adults in the phone-plan task, table~4 gives limited satisficing
at 70.7 percent under low accountability and 42.5 percent under high
accountability, while extensive noncompensatory search is 26.8 and
50.0 percent. Their distribution across strategy categories differed with
accountability, \(p=.04\). Across the tasks, extensive noncompensatory
search and compensatory search had similarly complex justifications,
and both exceeded limited satisficing. Two judges, blind to age and
accountability, scored differentiation and integration on a nine-point
scale; reported interjudge correlations were .94 and .82. Education,
not age, was the sole significant contributor in the authors' stepwise
strategy-classification analysis. That selected association does not
establish a causal effect of education
\cite[pp.~707--709, tables~3--5]{HessQueenEnnis2013Search}.

Crucially, participants were told there was no single correct choice,
and the study did not assess decision quality. Its strongest direct result
concerns sensitivity of search and explanation to context. Resource
conservation and adaptive use of preferences are interpretations, not
measured improvements in welfare or proof of a universal shift in goals
within each aging individual. Changing goals remain a possible mechanism
of observed performance, alongside cognitive resources and knowledge.

\section{When additional thought helps: information structure and task load}
\label{sec:p19-aging-task-structure}

Hess, Queen and Patterson tested 163 community-residing adults: 54 ages
20--44, 55 ages 45--64 and 54 ages 65--85. Education and vocabulary did
not significantly differ by age; working-memory, speed and executive-test
performance did. Random assignment within age group compared three
minutes of recorded think-aloud deliberation with three minutes of word
scrambles intended to prevent further deliberate processing of the choice.
The comparison does not directly observe an unconscious reasoning
mechanism \cite[Method]{HessQueenPatterson2012Deliberation}.

Each person considered a grocery-store choice and an apartment choice,
with task order and condition assignment counterbalanced. Four alternatives
were shown successively for twenty seconds each. Each had ten attributes,
five on common, alignable dimensions and five unalignable features. In
one task condition the informative, important common dimensions supported
the experimenter's best option. In the other, these common dimensions did
not distinguish the alternatives, so the less directly comparable features
had to be considered. After the intervening activity, participants chose
one option and rated all four without seeing their descriptions again.
The benchmark is the authors' constructed evaluative ordering, not a
unique best apartment or store for every actual consumer.

\begin{hbexample}{def:p19-aging-alignable-case}{A shared dimension can help or fail to discriminate}
In the paper's alignable-focus materials, the four alternatives have,
respectively, four, three, two and one positive alignable features out
of five. Their total positive-feature counts over all ten features are
instead five, eight, four and six. Thus, counting every positive statement
favors a different alternative from attending to the important common
dimensions. In the unalignable-focus materials, all four alternatives
have four positive alignable features; these features no longer identify
a winner. The distinction concerns \emph{which information discriminates
under the declared benchmark}, not merely whether a person thinks for
longer \cite[table~2]{HessQueenPatterson2012Deliberation}.
\end{hbexample}

Pooled across the age groups, optimal-choice proportions were .52 without
additional deliberation and .46 with deliberation in the alignable-focus
condition. The corresponding unalignable-focus proportions were .34 and
.53. The reported logistic comparison favored deliberation in the latter
condition, odds ratio 2.17, \(p=.02\), but not in the former,
\(p=.48\). The paper does not give a confidence interval for that odds
ratio. No unique age effect appeared when age was added to this choice
model. These results do not establish equality of age populations; they
show that a broad older-adult deficit was not detected in these particular
choices, despite age differences on the separate resource measures
\cite[Results, table~3]{HessQueenPatterson2012Deliberation}.

Rating analyses revealed more conditional differences: in the
unalignable-focus/no-deliberation condition, older adults with lower
education discriminated the alternatives less systematically; in the
alignable-focus/active-deliberation condition, older adults with higher
measured ability showed stronger discrimination. Education and ability
were observed characteristics, not assigned interventions. Analyses that
divided active deliberators according to their spoken content were
exploratory median-split analyses, not randomized comparisons of effective
and ineffective strategies. They suggest that attention to relevant
features matters, but cannot isolate the causal benefit of mentioning a
particular feature.

The same experiment also measured source memory after the decisions.
Choice-supportive memory for alignable features showed an age-related
pattern that was no longer significant after controlling for ability.
Neither that statistical attenuation nor successful choice identifies
emotion-regulation goals as the unique cause. Good immediate evaluation
can coexist with less accurate later attribution of features. Familiar
contexts, informative common dimensions and immediate impression
formation are plausible supports here; the study does not promise that
removing time pressure or instructing more thought improves every task.
In particular, distraction was not generally superior to deliberate
comparison \cite[Results and Discussion]{HessQueenPatterson2012Deliberation}.

\section{Accumulated expertise: baseline performance is not a rate of aging}
\label{sec:p19-aging-expertise}

Taylor, Kennedy, Noda and Yesavage followed general aviation pilots rather
than randomly selected older adults. Entry required ages 40--69, current
medical certification, current flying and 300--15,000 total flight hours;
retired major-airline pilots were excluded. Of 141 people completing
baseline testing before the specified cutoff, 118 had at least three
annual observations and entered the longitudinal analysis. Their mean
follow-up was 3.1 years, with 3.8 observations. The expertise strata were
32 pilots with basic visual-flight ratings, 60 with instrument ratings,
and 26 with instructor or air-transport certification. Groups differed in
age and flight experience; their mean education was approximately sixteen
to seventeen years, and most were white men. This selected, actively
practicing sample is not evidence about all adults or the oldest old
\cite[pp.~648--650, table~1]{TaylorKennedyNodaYesavage2007Pilot}.

After six familiarization flights, each annual session included morning
and afternoon seventy-five-minute simulator flights. The scenario combined
air-traffic instructions, traffic avoidance, engine emergencies and landing
under substantial workload. Different versions reduced memorization of
particular messages. Twenty-three recorded variables were standardized
using baseline means and standard deviations, aggregated into four
components and then averaged for the primary summary score. This is
specialized coordination of information and actions, not a test of
formal deduction alone.

The authors estimated each pilot's baseline level and annual slope, then
modeled these separately. In their notation, \(\beta_1\) is the
age coefficient and \(\beta_2\) the coefficient of ordinal expertise;
age was centered at the median and expertise coded as minus one, zero
and one. For the summary score, older age predicted lower baseline
performance, \(\beta_1=-.038\), standard error .006,
\(p<.0001\), but \emph{less} decline over the observed follow-up,
\(\beta_1=.004\), standard error .002, \(p<.01\).
Higher expertise predicted a better baseline, \(\beta_2=.155\),
standard error .060, and less decline, \(\beta_2=.039\), standard
error .015; both reported \(p<.05\). The signs refer to different
outcomes: a negative baseline age coefficient cannot be read as the
within-person annual rate of decline
\cite[pp.~650--651, table~2]{TaylorKennedyNodaYesavage2007Pilot}.

Components also differed. Communication accuracy declined most strongly,
while traffic avoidance improved, particularly among older pilots.
Replacing expertise ratings with flight-hour quantities did not reproduce
the significant expertise prediction. That comparison suggests that
activity amount and structured expertise are not interchangeable, but
does not randomize advanced training or identify its causal mechanism.
The age-by-expertise interaction was not detected, \(p=.15\).
Thus, the direct evidence supports an expertise advantage and different
trajectories, not a demonstrated greater protective effect of expertise
for older than for younger pilots.

The authors discuss hardy-survivor selection, cohort differences,
restricted ages, practice, floor effects and regression to the mean.
Failure to detect selective attrition among the 23 excluded participants
does not eliminate selection concerns. Traffic avoidance had relatively
low consistency, making regression to the mean a particular concern.
Continuing flight activity and its type were not assigned. Improved
traffic avoidance may reflect learning the repeated simulator rather than
a reversal of general cognitive aging. The small cells at the age and
expertise extremes, sometimes five to nine people, limit interaction
power. Compensation through knowledge or better strategies is a plausible
account requiring further discrimination, not the only explanation of
the fitted slopes \cite[pp.~652--654]{TaylorKennedyNodaYesavage2007Pilot}.

\section{A randomized training result and its transfer boundary}
\label{sec:p19-aging-training}

The ACTIVE ten-year report concerns independently living volunteers aged
at least 65 in six United States metropolitan areas. Recruitment occurred
in 1998--1999. The initial randomized total was 2,832; the flow diagram
identifies thirty protocol violations, leaving the 2,802-person analytical
sample. Its four arms contained 703 memory-training, 699 reasoning-training,
702 processing-speed-training and 698 no-contact-control participants.
Mean baseline age was about 73.6; approximately 76 percent were women
and 26 percent African American. Exclusions included substantial
baseline cognitive or functional impairment and conditions or sensory
limitations interfering with the protocol. These criteria select the
trial's population; they are not diagnostic rules supplied to this reader
\cite[Methods, Results, figure~1 and table~1]{RebokEtAl2014ACTIVE}.

Initial training comprised ten small-group sessions of sixty to
seventy-five minutes over five to six weeks. Reasoning training addressed
problems with serial patterns. Its outcome combined total-correct scores
on Letter Series, Letter Sets and Word Series; memory and processing
speed had different tasks. Four-session boosters at eleven and
thirty-five months were offered to a randomly selected subset of people
who had completed at least eight initial sessions. Eligibility for this
second selection depended on initial completion.

Repeated-measures mixed models used the assigned groups, baseline
covariates and available observations, with multiple-imputation and
low-missing-score sensitivity analyses. The reported training effect size
is training-group improvement from baseline to year ten minus
control-group improvement, divided by the within-person standard
deviation of the composite. It is a \emph{between-group contrast of
change}, not an assertion that every trained person improved by that
amount. The authors used 99 percent confidence intervals and a
multiplicity-adjusted significance criterion
\cite[Methods and table~2]{RebokEtAl2014ACTIVE}.

At ten years the reasoning-training contrast on the reasoning outcome
was .23, with 99 percent confidence interval .09--.38. Table~2 gives
938 complete cases for that outcome \emph{across groups}, not 938
reasoning-trained participants. The corresponding targeted memory
contrast was .06, interval -.14--.27, and the targeted speed contrast
was .66, interval .43--.88. These are different measured outcomes, not
three estimates of a general reasoning improvement. The reported
reasoning-group mean change was approximately flat while the control
declined; a retained relative advantage need not be an absolute increase
from baseline.

The reasoning group also reported less difficulty with instrumental
daily activities: effect-size contrast .38, 99 percent interval .02--.74.
But the performance-based everyday problem-solving contrast was -.02,
interval -.25--.22, and the study reported no training effect on its
performance-based everyday-function measures. Self-reported function,
tested serial-pattern skill and observed performance on other tasks
therefore retain distinct transfer standings. The exploratory future
health implications discussed by the authors are not measured prevention
of dementia in this report \cite[table~2 and Discussion]{RebokEtAl2014ACTIVE}.

Ten-year retention was 44 percent. Death, withdrawal and continuing
missed visits were major causes of nonparticipation; retained people
were initially younger and healthier. Similar attrition across arms and
the reported sensitivity analyses provide useful support for the trial
contrast, but do not guarantee generalization to excluded or missing
people. The booster comparison has an additional qualification: the
nonbooster group includes many initial noncompleters, so the compared
groups do not have the same initial training exposure. The authors
explicitly caution against an unqualified booster interpretation.
Finally, ACTIVE has no younger-adult comparison group. It establishes
targeted training effects in its older-adult sample, not elimination of
an older--younger performance gap or a universally effective intervention.

\section{Using the evidence without flattening its differences}
\label{sec:p19-aging-use}

\begin{hbmethod}{def:p19-aging-lookup}{Evaluate an age-sensitive reasoning support}
Specify the reasoning or decision criterion, the age range and recruitment,
and the task's time, memory, interface and knowledge demands. Separate
changing \emph{willingness} to engage from changing available processing
resources. Compare the same outcome under the relevant support, preserving
within-person and repeated-item dependence. To claim reduced age
differences, examine the age-by-condition comparison; improvement within
older adults alone is insufficient. Distinguish a nonsignificant difference
from evidence of equivalence. For training, state randomization, the
comparison activity, completion, retention, duration, uncertainty and
whether transfer was actually tested. For expertise, retain the possibility
of selection, prior ability and practice rather than treating experience
as a randomized treatment.
\end{hbmethod}

These records suggest specific conditions under which a blanket decline
claim is a poor description: personally relevant decisions can elicit
more engaged search; informative common dimensions can support immediate
comparison; practiced domain knowledge can accompany preserved specialized
performance; and targeted training can preserve an advantage on the
trained skill. Their failure boundaries are equally important. More
search need not improve choice quality. Additional thought about
irrelevant information need not help. Expertise need not moderate an
age effect, and an intervention's near-domain benefit need not transfer
to every everyday task. None makes chronological age a clinical diagnosis
or establishes wisdom from a favorable score.

\begin{proofboundary}{def:p19-aging-evidence-boundary}{Four primary records, distinct conclusions}
The search study directly measures engagement and explanations, not optimal
choices. The deliberation study measures constructed choice and rating
criteria, with relatively small detected age effects and conditional
education/ability associations. The simulator study measures longitudinal
specialized performance in a selected cohort, without randomized expertise
or a detected age-by-expertise interaction. ACTIVE supplies randomized
training contrasts, with attrition and domain-specific transfer limits,
but no younger comparison. These outcomes neither establish a universal
aging trajectory nor supply individual medical or occupational advice.
No mathematical theorem about human cognition, participant-data reanalysis
or new replication is claimed.
\end{proofboundary}

\index{aging!goals and reasoning}
\index{cognitive compensation}
\index{expertise!aging}
\index{reasoning training!older adults}
